% Generated draft: requires mathematical and source review.
\documentclass[11pt]{article}
\usepackage[margin=25mm]{geometry}
\usepackage{fontspec}
\setmainfont{Latin Modern Roman}
\usepackage{amsmath,amssymb,mathtools,unicode-math}
\setmathfont{Latin Modern Math}
\usepackage{enumitem,needspace,xurl,hyperref}
\setlength{\parindent}{0pt}
\setlength{\parskip}{5pt}
\setlength{\emergencystretch}{3em}
\newcommand{\N}{\mathbb{N}}
\newcommand{\Z}{\mathbb{Z}}
\newcommand{\Q}{\mathbb{Q}}
\newcommand{\R}{\mathbb{R}}
\newcommand{\C}{\mathbb{C}}
\newcommand{\F}{\mathbb{F}}
\newcommand{\A}{\mathbb{A}}
\newcommand{\PP}{\mathbb{P}}
\newcommand{\E}{\mathbb{E}}
\newcommand{\eps}{\varepsilon}
\newcommand{\dd}{\,\mathrm{d}}
\begin{document}
\section*{UnsolvedMath ID: 30003493}
\textbf{Source problem code:} catalogue entry
\subsection{Definitions and mathematical statement}
Write
\[
A=\C[G/U].
\]
The supplied record denotes by this expression the algebra associated with the quotient \(G/U\), but does not specify the hypotheses on \(G\), the subgroup \(U\), or the precise quotient and coordinate-ring conventions. These data must be recovered from the source before a fully specified formulation is possible.

Let \(B_{\mathrm{can}}\) denote the basis of \(A\) described in the record as constructed by cluster duality, and let \(B\) denote the basis described there as the classical canonical basis of \(A\). A \(\C\)-basis here means a subset such that every element of \(A\) has a unique expression as a finite \(\C\)-linear combination of its elements. The particular constructions and normalizations identifying these two named bases are not supplied.

\textbf{Problem (as supplied).} Determine the relationship between \(B_{\mathrm{can}}\) and \(B\) in \(\C[G/U]\).

The source poses a comparison question; it does not explicitly conjecture \(B_{\mathrm{can}}=B\), specify another mathematical relation to establish, or quantify over a stated class of groups. Consequently, equality cannot be substituted for the requested relationship, and no universal quantifier over groups can be added without further source information. The exact classical canonical-basis convention, including whether a dual canonical basis is intended for this function algebra, also requires clarification.
\subsection{Short English statement}
The problem asks how the basis of \(\C[G/U]\) obtained through cluster duality relates to the classical canonical basis named in the source. The group hypotheses, the exact basis constructions, and the intended comparison are unspecified, so this remains a comparison question requiring source clarification and human mathematical review.
\subsection{Sources}
\begin{itemize}
\item Local UnsolvedMath registry: supplied complete record for permanent ID \(30003493\), title \emph{Cluster Duality and Classical Canonical Bases}, source problem code \texttt{catalogue entry}, in the Oberwolfach Reports Open Problems collection. No individual proposer is identified; the supplied proposed year is 2017.
\item Supplied source citation: \emph{Algebraic Groups (2018)}, Oberwolfach Reports, pp.\ 1281--1347, DOI \texttt{10.4171/owr/2017/21}. \url{https://doi.org/10.4171/owr/2017/21}. The literature-source list also labels this reference ``Oberwolfach Reports 2018.'' The supplied 2018 label, 2017 DOI component, and proposed year are retained without resolving their bibliographic relationship.
\item Supplied literature reference: Genz et al., ``canonical bases and cluster duality.'' \url{https://arxiv.org/abs/1711.07176}. This is the supplied attribution and descriptive label; full bibliographic details are not provided.
\item Supplied literature reference: Genz et al., ``type A canonical bases.'' \url{https://arxiv.org/abs/1611.03465}. This is the supplied attribution and descriptive label; full bibliographic details are not provided.
\item Supplied literature assessment, checked \textbf{2026-08-21}: the record labels the problem open and reports that later work connects potential functions with crystal and Lusztig parametrizations, including type-specific results, without establishing the exact equality or general relationship asked about here. This is supplied dated curation derived from the Oberwolfach Reports collection, not an independently verified literature assessment. Its time-sensitive status and cited claims require independent verification.
\end{itemize}
\end{document}
